% The Turtle Book: a longer, guided introduction to Logo (with a little
% BASIC at the end) for a strong young reader. Companion to kidbook.tex and
% set with the same kidstyle.sty.
%
% Every program listed here is a file under programs/, and every picture
% under figures/ was drawn by running that file. To re-test the programs
% and regenerate the figures:  ../tools/check-programs.sh  (or make check).
%
% The grey "To write" boxes are the outline for the prose. Set
% \draftnotesfalse below to print the book without them. kidbook opens each
% chapter with \lede{...} for its first line; do the same here once written.

\documentclass[12pt,oneside,paper=5.5in:8.5in,DIV=11,chapterprefix=false,headings=big]{scrbook}

\usepackage{kidstyle}
\usepackage{graphicx}
\usepackage{booktabs}
\usepackage{array}
\hypersetup{pdftitle={The Turtle Book}}
% Twenty chapters: two-digit numbers need more room in the contents than
% kidstyle allows for.
\DeclareTOCStyleEntry[numwidth=2.3em]{chapter}{chapter}

% ---- Draft notes ----------------------------------------------------------
\newif\ifdraftnotes
\draftnotestrue
\newtcolorbox{draftbox}{%
  enhanced, breakable,
  colback=kidink!5, colframe=kidink!35, boxrule=0.6pt, arc=3pt,
  left=6pt, right=6pt, top=4pt, bottom=4pt,
  before skip=0.6em, after skip=0.6em,
  fontupper=\small, before upper=\RaggedRight}
\newcommand{\draftnote}[1]{\ifdraftnotes\begin{draftbox}\textbf{\color{kidink!70}To write:} #1\end{draftbox}\fi}

% ---- Programs and pictures ------------------------------------------------
% A listing is read straight from programs/, so the book shows exactly what
% tools/check-programs.sh ran. \small rather than kidbook's \normalsize: some
% BASIC lines here run to 52 characters.
\newcommand{\logoprogram}[1]{%
  \begin{kidcodebox}\VerbatimInput[fontsize=\small]{programs/logo/#1.lg}\end{kidcodebox}}
\newcommand{\basicprogram}[1]{%
  \begin{kidcodebox}\VerbatimInput[fontsize=\small]{programs/basic/#1.BAS}\end{kidcodebox}}
% What a Logo program drew. Optional argument caps the height. No extension:
% most figures are vector PDFs from the EPS export, but the few that need a
% real screenshot (fill, label) are PNGs, and graphicx finds either.
\newcommand{\drawing}[2][1.5in]{%
  \begin{center}
    \includegraphics[height=#1,width=0.8\linewidth,keepaspectratio]{figures/#2}
  \end{center}}
% What a program printed: a dark box, like the screen.
\newtcolorbox{screenbox}{%
  enhanced, colback=kidink!92, colframe=kidink!92, arc=5pt,
  left=10pt, right=10pt, top=6pt, bottom=6pt,
  before skip=0.9em, after skip=0.9em}
\newcommand{\screen}[1]{%
  \begin{screenbox}\VerbatimInput[fontsize=\small,formatcom=\color{white}]{figures/#1.txt}\end{screenbox}}

\begin{document}

\frontmatter

% ---------------------------------------------------------------------------
% Cover
% ---------------------------------------------------------------------------
\begin{titlepage}
\centering

\vspace*{\stretch{1}}

\begin{tikzpicture}
  \foreach \i/\c in {0/kidred, 1/kidorange, 2/kidgreen, 3/kidblue, 4/kidpurple, 5/kidteal}
    \fill[\c] (\i*0.62,0) circle (0.17);
\end{tikzpicture}

\vspace{1.2em}

{\fontsize{34}{38}\selectfont\bfseries\color{kidink} The Turtle\\[0.1em]
 \color{kidgreen} Book\par}

\vspace{0.8em}

{\large\bfseries\color{kidink!70} A longer book about Logo,\\
 with a little BASIC at the end\par}

\vspace{1.6em}

% The turtle part way round a square, the first shape in the book.
\begin{tikzpicture}[scale=1.0,line join=round]
  \draw[kidblue,line width=2.2pt,line cap=round]
    (-1.5,-1.5) -- (-1.5,1.5) -- (1.5,1.5) -- (1.5,-0.2);
  \begin{scope}[shift={(1.5,-0.55)},rotate=180,scale=0.8]
    \kidturtlebody{kidgreen}{kidgreendark}{kidgreendark}{white}
  \end{scope}
\end{tikzpicture}

\vspace{1.6em}

{\Large\bfseries\color{kidink} For \rule{5em}{0.5pt}\par}

\vspace*{\stretch{1.15}}
\end{titlepage}

\draftnote{Her name goes on the line on the cover. Decide whether this
  book is hers alone or shared; the kidbook names both girls.}

% ---------------------------------------------------------------------------

\renewcommand*{\contentsname}{What Is Inside}
\chaptericon{\faListUl}
\tableofcontents

\mainmatter

\part{The Turtle}

% ---------------------------------------------------------------------------
\chaptericon{\faPowerOff}
\chapter{Before You Start}

\draftnote{How to open Logo from the menu (``Draw Pictures with Logo
  Turtle''). What appears: a grey screen, a small triangle in the middle
  (that is the turtle, seen from above), a house it has already drawn,
  and a question mark where you type. You type a line, press
  \key{Enter}, the turtle does it. Mistakes are fine and expected; the
  worst that happens is a message (Chapter~\ref{ch:wrong}). How to
  leave: type \texttt{bye}, or hold \key{Ctrl} and \key{Alt} and press
  \key{Backspace}. \par
  Decide about case: this book prints commands in lowercase because it is
  easier to type, but the kidbook uses \texttt{FD 50}. Logo accepts both.
  Say so once.}

% ---------------------------------------------------------------------------
\chaptericon{\kidturtlew[0.72]}
\chapter{Meet the Turtle}

\draftnote{The turtle knows four moving words: \texttt{fd} (forward),
  \texttt{bk} (back), \texttt{rt} (right), \texttt{lt} (left). Each one
  wants a number after it: how far, or how much to turn. The turtle always
  has a place and a direction it is facing; \texttt{fd} changes the place,
  \texttt{rt} and \texttt{lt} change the direction. A good first activity
  away from the computer: one person is the turtle, the other gives
  commands, and the turtle may only do exactly what was said.}

Type these three lines, one at a time, pressing \key{Enter} after each:

\logoprogram{first-steps}
\drawing{first-steps}

\draftnote{Two more words: \texttt{cs} wipes the screen and puts the turtle
  back in the middle facing up; \texttt{home} sends it back without
  wiping.}

\begin{trythis}[Try]
\begin{itemize}
\item \chip{bk 100} --- does the turtle draw when it walks backwards?
\item \chip{rt 90} four times in a row. Where is it facing now?
\item Draw the letter L. Then the letter T (you will need \chip{bk}).
\end{itemize}
\end{trythis}

% ---------------------------------------------------------------------------
\chaptericon{\faRedo}
\chapter{Turning}

\draftnote{What the number after \texttt{rt} means. \texttt{rt 90} is a
  quarter turn, the corner of a square; \texttt{rt 180} is turning to face
  the opposite way; \texttt{rt 360} is all the way around, so the turtle
  ends up facing where it started. Smaller numbers are smaller turns.
  \texttt{rt 45} is half of a corner. Do it with your body: face the
  window, \texttt{rt 90}, now you face the door.}

\logoprogram{zigzag}
\drawing{zigzag}

\begin{trythis}[Try]
\begin{itemize}
\item Change every \chip{50} to \chip{30}. Then to \chip{80}.
\item Change \chip{45} to \chip{20}. What changed?
\item Make the zigzag longer by typing more lines.
\end{itemize}
\end{trythis}

% ---------------------------------------------------------------------------
\chaptericon{\faIcon{sync-alt}}
\chapter{Saying It Again}

Here is a square, the long way:

\logoprogram{square-long}

\draftnote{Notice the same two lines typed four times. Logo has a word for
  ``do this again'': \texttt{repeat}, then how many times, then the
  commands inside square brackets. The brackets say where the repeated
  part starts and ends.}

\logoprogram{square-repeat}
\drawing{square-repeat}

\draftnote{The big idea in this chapter: to draw a closed shape and end up
  facing the way it started, the turtle turns all the way around once,
  which is 360. Four corners share 360, so each is 90. Three corners share
  360, so each is 120. Let her check the table with a calculator or by
  counting up.}

\logoprogram{triangle}
\drawing[1.2in]{triangle}

\logoprogram{hexagon}
\drawing[1.2in]{hexagon}

\begin{center}
\begin{tabular}{rrl}
\toprule
sides & turn & \\
\midrule
3 & 120 & triangle \\
4 & 90 & square \\
5 & 72 & pentagon \\
6 & 60 & hexagon \\
8 & 45 & octagon \\
12 & 30 & \\
\bottomrule
\end{tabular}
\end{center}

\draftnote{Sides times turn is always 360. Invite her to guess what
  twelve sides will look like before running it.}

\logoprogram{twelve-sides}
\drawing[1.2in]{twelve-sides}

\draftnote{And if the turtle takes 360 tiny steps with 360 tiny turns, the
  corners disappear.}

\logoprogram{circle}
\drawing[1.2in]{circle}

\begin{trythis}[Try]
\begin{itemize}
\item \chip{repeat 4 [fd 100 rt 90]} but change \chip{90} to
  \chip{91}. Then \chip{repeat 40} of that.
\item \chip{repeat 5 [fd 100 rt 144]} --- a star. Why 144 and not 72?
\item \chip{repeat 360 [fd 2 rt 1]} --- a bigger circle.
\end{itemize}
\end{trythis}

% ---------------------------------------------------------------------------
\chaptericon{\faPaintBrush}
\chapter{The Pen}

\draftnote{The turtle carries a pen. \texttt{pu} lifts it so the turtle
  moves without drawing; \texttt{pd} puts it back down. This is how you
  draw two things that do not touch.}

\logoprogram{dashes}
\drawing[1.2in]{dashes}

\draftnote{\texttt{setpc} sets the pen colour by number and
  \texttt{setpensize} makes the line thicker. The two numbers in
  \texttt{setpensize [5 5]} are width and height; just keep them the
  same.}

\logoprogram{colors}
\drawing[1.3in]{colors}

\begin{center}
\begin{tabular}{rlrl}
\toprule
0 & black   & 8  & brown \\
1 & blue    & 9  & tan \\
2 & green   & 10 & forest green \\
3 & light blue & 11 & aqua \\
4 & red     & 12 & salmon \\
5 & magenta & 13 & purple \\
6 & yellow  & 14 & orange \\
7 & white   & 15 & grey \\
\bottomrule
\end{tabular}
\end{center}

\draftnote{\texttt{setbg} changes the background with the same numbers.
  The screen starts grey (15) with a black pen because the welcome file
  sets it that way.}

\draftnote{Inside a \texttt{repeat}, the word \texttt{repcount} is the
  number of the turn Logo is on: 1 the first time, 2 the second, and so
  on. Using it as the pen colour gives a different colour each spoke.}

\logoprogram{color-wheel}
\drawing[1.3in]{color-wheel}

\draftnote{\texttt{fill} pours colour into the closed shape the turtle is
  standing inside, so first draw the shape, then lift the pen and step
  inside. If the turtle is outside the shape, the colour goes outside.}

\logoprogram{fill}
\drawing[1.2in]{fill}

\begin{trythis}[Try]
\begin{itemize}
\item A dotted circle: \chip{repeat 36 [fd 5 pu fd 5 pd rt 10]}
\item Every colour: \chip{repeat 16 [setpc repcount fd 20 rt 90 fd 5 lt 90]}
\item Fill the house that was on the screen when Logo opened.
\end{itemize}
\end{trythis}

% ---------------------------------------------------------------------------
\chaptericon{\faGraduationCap}
\chapter{Teaching the Turtle a New Word}
\label{ch:to}

\draftnote{The turtle only knows the words it was born with, but you can
  teach it more. \texttt{to square} starts a lesson; the lines after it
  are what \texttt{square} means; \texttt{end} finishes the lesson. While
  you are teaching, the prompt changes from \texttt{?} to \texttt{>} so
  you know Logo is listening to the lesson rather than doing things. After
  \texttt{end}, \texttt{square} is a word like any other.}

\logoprogram{square-word}
\drawing[1.3in]{square-word}

\draftnote{Useful words about words: \texttt{po "square} shows the lesson
  back to you (the quote mark is part of it); \texttt{pots} lists every
  word you have taught; \texttt{erase "square} forgets one. \par
  Words are forgotten when Logo closes. \texttt{save "mine.lg} writes them
  all to a file and \texttt{load "mine.lg} brings them back next time.
  The file lands wherever Logo was started from, which the menu keeps
  the same every time, so \texttt{load} finds it again.}

\begin{trythis}[Try]
\begin{itemize}
\item Teach the turtle \chip{triangle}. Then \chip{hexagon}.
\item \chip{repeat 8 [square rt 45]}
\item Teach it \chip{box}: a square, then back to where it started,
  facing the same way. Check with \chip{repeat 3 [box fd 120]}.
\end{itemize}
\end{trythis}

% ---------------------------------------------------------------------------
\chaptericon{\faHashtag}
\chapter{Words That Take a Number}

\draftnote{\texttt{fd} takes a number; why can't \texttt{square}? It can.
  \texttt{to square :size} means ``square wants one number, and inside
  the lesson we will call it \texttt{size}.'' The colon (dots) means ``the
  number that was handed to me,'' not the word itself. Then
  \texttt{square 30} and \texttt{square 120} are the same lesson with
  different numbers.}

\logoprogram{square-size}
\drawing[1.4in]{square-size}

\draftnote{A word you taught can be used inside another lesson, exactly
  like a built-in word. \texttt{flower} hands its own number on to
  \texttt{square}.}

\logoprogram{flower}
\drawing[1.5in]{flower}

\begin{trythis}[Try]
\begin{itemize}
\item \chip{flower 30}, \chip{flower 120}
\item Change \chip{12} and \chip{30} in \chip{flower} (keep them
  multiplying to 360, or don't, and see).
\item Teach \chip{triangle :size} and make a triangle flower.
\end{itemize}
\end{trythis}

% ---------------------------------------------------------------------------
\chaptericon{\faHome}
\chapter{Building Big Things}

\draftnote{Big pictures are made from small words. A house is a square with
  a triangle on top; the tricky part is walking the turtle to the top-left
  corner and tilting it 30 degrees so the roof sits flat. Trace it with a
  finger, line by line, saying where the turtle is. This is the house that
  is already on the screen when Logo opens.}

\logoprogram{house}
\drawing[1.5in]{house}

\draftnote{\texttt{house} leaves the turtle exactly where it started,
  facing the same way. That is on purpose: it makes \texttt{house} safe to
  use inside a \texttt{repeat}. Walk right with the pen up, draw another.}

\logoprogram{street}
\drawing[1.2in]{street}

\begin{trythis}[Try]
\begin{itemize}
\item A street of five houses. (Do they fit? Make them smaller.)
\item A door: a small tall rectangle. Where does the turtle have to walk first?
\item A window, a chimney, a sun (\chip{circle}) in the sky.
\end{itemize}
\end{trythis}

% ---------------------------------------------------------------------------
\chaptericon{\faInfinity}
\chapter{A Word That Uses Itself}

\draftnote{The strangest and best trick. \texttt{spiral} draws one side,
  turns, and then calls \texttt{spiral} again with a slightly bigger
  number. Each square is a little bigger than the last, so the line
  spirals out. \par
  The \texttt{if} line is the stop rule: once the number gets past 150,
  \texttt{stop} ends the lesson and the turtle rests. Ask what would
  happen without it: it would never stop. The way out is the
  \textbf{Logo} menu at the top of the window, then \textbf{Stop}
  (verified on this version of UCBLogo; Ctrl+C does nothing while a
  program runs). The menu shows Alt+S as a shortcut, still to be tried
  on the Pi. If all else fails, \key{Ctrl}+\key{Alt}+\key{Backspace}
  leaves Logo for the menu, as the kidbook says.}

\logoprogram{spiral}
\drawing[1.5in]{spiral}

\logoprogram{triangle-spiral}
\drawing[1.5in]{triangle-spiral}

\begin{trythis}[Try]
\begin{itemize}
\item In \chip{spiral}, change \chip{rt 90} to \chip{rt 91}. Then \chip{rt 89}.
\item Change \chip{+ 10} to \chip{+ 3}.
\item Change \chip{150} to \chip{300}.
\item Make a \chip{hexspiral} with \chip{rt 60}.
\end{itemize}
\end{trythis}

% ---------------------------------------------------------------------------
\chaptericon{\faCommentDots}
\chapter{The Turtle Can Talk}

\draftnote{\texttt{print} writes in the text part of the screen. Words go
  in square brackets. Numbers can be added: \texttt{sum 2 3} or
  \texttt{2 + 3}, times is \texttt{*}. So Logo is also a calculator.}

\logoprogram{talking}
\screen{talking}

\draftnote{\texttt{label} writes words on the drawing, at the turtle,
  facing the way the turtle faces. \texttt{setlabelheight} makes the
  letters bigger. Put her name in it.}

\logoprogram{name-sign}
\drawing[1in]{name-sign}

\begin{trythis}[Try]
\begin{itemize}
\item \chip{print 100 * 100}
\item \chip{rt 90 label [sideways]}
\item Write your name at the four corners of a square:
  \chip{repeat 4 [label [ME] fd 100 rt 90]}
\end{itemize}
\end{trythis}

% ---------------------------------------------------------------------------
\chaptericon{\faDice}
\chapter{Surprises}

\draftnote{\texttt{random 30} gives a number the turtle picks by itself,
  from 0 up to 29 (not 30; say so, she will ask). Run \texttt{print random
  10} several times to see. A program with \texttt{random} in it draws a
  different picture every time.}

\logoprogram{wander}
\drawing[1.4in]{wander}

\logoprogram{random-colors}
\drawing[1.4in]{random-colors}

\begin{trythis}[Try]
\begin{itemize}
\item Run \chip{repeat 100 [wander]} three times. Same picture?
\item \chip{repeat 20 [square random 100]}
\item \chip{repeat 50 [setpc random 16 house random 80]}
\end{itemize}
\end{trythis}

% ---------------------------------------------------------------------------
\chaptericon{\faGamepad}
\chapter{Driving the Turtle}

\draftnote{A game. \texttt{readchar} waits for one key and hands it over;
  \texttt{make "key} stores it under the name \texttt{key}, and
  \texttt{:key} reads it back, the same colon as in \texttt{:size}. Each
  \texttt{if} checks for one letter. The last line calls \texttt{drive}
  again, so it keeps listening until you press \key{Q}. Keys: \key{W}
  forward, \key{S} back, \key{A} left, \key{D} right, \key{Q} quit.
  The picture is what the check script drove; hers will differ.}

\logoprogram{drive-interactive}
\drawing[1.3in]{drive-interactive}

\begin{trythis}[Try]
\begin{itemize}
\item Make \key{W} go further. Make \key{A} and \key{D} turn less.
\item Add a key that lifts the pen and one that puts it down.
\item Add a key that changes colour.
\item Add \key{C} for \chip{cs}.
\end{itemize}
\end{trythis}

% ---------------------------------------------------------------------------
\chaptericon{\faExclamationTriangle}
\chapter{When Things Go Wrong}
\label{ch:wrong}

\draftnote{The turtle never breaks; it just says it does not understand.
  These are the four messages she will see most. Each one is the turtle
  being honest about exactly what confused it.}

\begin{center}
\small
\begin{tabular}{>{\raggedright\arraybackslash}p{1.55in}>{\raggedright\arraybackslash}p{2.35in}}
\toprule
The turtle says & What it means \\
\midrule
\texttt{I don't know how to squar}
  & A word it has not learned. Usually a spelling slip, or a lesson you
    have not taught yet in this session. \\[0.5em]
\texttt{not enough inputs to fd}
  & \texttt{fd} wanted a number and did not get one. \\[0.5em]
\texttt{You don't say what to do with 50}
  & A number with no word in front of it. The turtle has 50 and no
    idea what for. \\[0.5em]
\texttt{size has no value}
  & \texttt{:size} used outside a lesson that was handed a size. \\
\bottomrule
\end{tabular}
\end{center}

\draftnote{Also: if the turtle will not stop, use the Logo menu at the
  top, then Stop. A missing \texttt{]} makes Logo wait silently for more
  (the prompt changes); type \texttt{]} and Enter. And the difference
  between \texttt{square} and \texttt{"square}: without the quote, Logo
  runs the word; with it, Logo means the name.}

% ---------------------------------------------------------------------------
\chaptericon{\faBook}
\chapter{Words the Turtle Knows}

\section*{Moving}
\begin{commandlist}
\cmd{fd 50}{forward}
\cmd{bk 50}{back}
\cmd{rt 90}{turn right}
\cmd{lt 90}{turn left}
\cmd{cs}{clear the screen, turtle to the middle}
\cmd{home}{turtle to the middle, keep the drawing}
\cmd{ht}{hide the turtle}
\cmd{st}{show the turtle}
\end{commandlist}

\section*{The pen}
\begin{commandlist}
\cmd{pu}{pen up}
\cmd{pd}{pen down}
\cmd{setpc 4}{pen colour}
\cmd{setbg 7}{background colour}
\cmd{setpensize [5 5]}{thicker pen}
\cmd{fill}{pour colour into the shape you are inside}
\end{commandlist}

\section*{Again and again}
\begin{commandlist}
\cmd{repeat 4 [ ... ]}{do it again}
\cmd{repcount}{which time this is, inside a repeat}
\cmd{to name :size}{start teaching a word}
\cmd{end}{finish teaching}
\cmd{po "name}{show a lesson}
\cmd{pots}{list all lessons}
\cmd{erase "name}{forget a lesson}
\cmd{save "f.lg}{keep lessons in a file}
\cmd{load "f.lg}{get them back}
\cmd{if :size > 150 [stop]}{stop rule}
\end{commandlist}

\section*{Talking}
\begin{commandlist}
\cmd{print [hi]}{write words}
\cmd{label [hi]}{write words on the drawing}
\cmd{random 30}{a number from 0 to 29}
\cmd{make "key readchar}{wait for a key and remember it}
\cmd{wait 60}{pause one second}
\cmd{bye}{leave Logo}
\end{commandlist}

\part{A Little BASIC}

% ---------------------------------------------------------------------------
\chaptericon{\faTerminal}
\chapter{A Different Language}

\draftnote{BASIC is another computer language on the same computer (``Write
  BASIC Programs'' on the menu). It is better at words and questions than
  at drawing. Two differences from Logo up front: every line of a program
  starts with a number, and the computer does nothing until you type
  \texttt{RUN}. The numbers say what order to go in; counting by tens
  leaves room to squeeze a line in between later. \texttt{LIST} shows the
  program, \texttt{NEW} throws it away. \par
  Stopping a runaway program: \key{Ctrl}+\key{C}, as the kidbook says
  (PC-BASIC treats it like Ctrl+Break by default, so the welcome
  screen's Ctrl+Break works too).}

\basicprogram{HELLO}
\screen{basic-HELLO}

\draftnote{Line 40: with quotes, \texttt{PRINT} shows the letters; without
  quotes, it works out the sum first. Same trick as Logo's
  \texttt{print 2 + 3}.}

\begin{trythis}[Try]
\begin{itemize}
\item Add a line 45 that prints your age. Type \chip{LIST}. Where did it go?
\item \chip{PRINT 100 * 100} with no line number. It runs right away.
\item Type line 20 again with different words. Which one wins?
\end{itemize}
\end{trythis}

% ---------------------------------------------------------------------------
\chaptericon{\faQuestionCircle}
\chapter{The Computer Asks}

\draftnote{\texttt{INPUT} prints the question, waits for an answer, and
  stores it in \texttt{N\$}. The dollar sign means ``this box holds
  words''; a box for numbers has no dollar sign. Semicolons in
  \texttt{PRINT} glue pieces together on one line.}

\basicprogram{NAME}
\screen{basic-NAME}

\begin{trythis}[Try]
\begin{itemize}
\item Ask for a favourite animal too, and print both.
\item Print the name ten times (the kidbook's \chip{FOR} loop).
\end{itemize}
\end{trythis}

% ---------------------------------------------------------------------------
\chaptericon{\faPenFancy}
\chapter{Mad Libs}

\draftnote{Four questions, then a story with the answers dropped in. This
  is the program she will want to change most, so leave room. Note line
  80: \texttt{N} is a number box, so no dollar sign, and BASIC puts a space
  on each side of a printed number by itself.}

\basicprogram{MADLIB}
\screen{basic-MADLIB}

\begin{trythis}[Try]
\begin{itemize}
\item Write your own story with five blanks.
\item Ask for a name and make the story about that person.
\end{itemize}
\end{trythis}

% ---------------------------------------------------------------------------
\chaptericon{\faBullseye}
\chapter{Guess My Number}

\draftnote{Three new pieces: \texttt{RND} picks the secret number (line 30
  turns a fraction into 1 to 10); \texttt{IF \ldots\ THEN} does the rest
  of the line only when the test is true; \texttt{GOTO 50} jumps back to
  the question. Ask why the program needs \texttt{GOTO} here (it must ask
  again) and why line 80 does not (it is finished).}

\basicprogram{GUESS}
\screen{basic-GUESS}

\begin{trythis}[Try]
\begin{itemize}
\item Make it 1 to 100. How many guesses do you need if you are clever?
\item Count the guesses: add \chip{T = T + 1} and print \chip{T} at the end.
\item Swap roles: you pick, the computer guesses (hard; do it together).
\end{itemize}
\end{trythis}

% ---------------------------------------------------------------------------
\chaptericon{\faMusic}
\chapter{Counting Down, Colours, Music}

\draftnote{\texttt{FOR I = 10 TO 1 STEP -1} counts backwards. Line 40
  is a musical rest: \texttt{PLAY "P2"} plays silence for half a note,
  which is one second, and that is how the countdown gets its beat
  (a do-nothing loop is no good; this computer runs it in no time).
  \texttt{BEEP} at the end is the speaker. Whether the beep is audible
  depends on the Pi's sound, see doc/PI-CHECKLIST.txt.}

\basicprogram{COUNTDOWN}

\draftnote{\texttt{COLOR} changes the colour of the text that comes after
  it, numbers 1 to 15 like Logo's pen. \texttt{COLOR 7} puts it back to
  normal.}

\basicprogram{COLORS}

\draftnote{\texttt{PLAY} takes a string of note names. A number after a
  note makes it longer (2 is a half note). This one is Twinkle Twinkle.
  \textbf{The one thing left to check on the Pi}: that PC-BASIC actually
  makes sound there (doc/PI-CHECKLIST.txt). If it does not, drop this
  section and the \texttt{BEEP} in the countdown.}

\basicprogram{MUSIC}

\begin{trythis}[Try]
\begin{itemize}
\item Count down from 20. Count up.
\item Print your name in every colour.
\item \chip{PLAY "C D E F G A B > C"} --- a scale. What does \chip{>} do?
\item Write the next line of Twinkle Twinkle.
\end{itemize}
\end{trythis}

% ---------------------------------------------------------------------------
\chaptericon{\faBalanceScale}
\chapter{Two Languages, Same Ideas}

\draftnote{A closing page. Both languages have a way to say ``again''
  (\texttt{repeat} / \texttt{FOR}), a way to name a thing
  (\texttt{make} / \texttt{N = }), a way to choose (\texttt{if} /
  \texttt{IF THEN}), and a way to ask (\texttt{readchar} /
  \texttt{INPUT}). The words differ; the ideas are the same, and they are
  the same in the languages Daddy uses at work. What to do next: pick one
  of the ``Try'' ideas and make it yours.}

% ---------------------------------------------------------------------------
% Write-in pages, as in kidbook. pdfjam pads the rest to a multiple of 4.

\dotgridpage{My Turtle Words}
\dotgridpage{My Programs}

% Colophon
\thispagestyle{empty}
\null\vfill
{\centering\footnotesize\color{kidink!55}
 The Turtle Book\\
 Updated \today\par}

\end{document}
